\documentclass[11pt]{article}
\usepackage[margin=1in]{geometry}
\usepackage{booktabs}
\usepackage{graphicx}
\usepackage{hyperref}

\title{RoofSight: Ground-View Roof and Rooftop-Obstacle Segmentation for On-Site PV Planning}
\author{}
\date{}

\begin{document}
\maketitle

\begin{abstract}
Draft. A dataset of ground-level photographs of residential roofs with instance masks for roof
planes, rooftop obstacles, vegetation occlusion and roof edges; a benchmark with fixed splits;
baselines with RF-DETR-Seg, SAM~3 zero-shot, YOLO26-seg and Mask2Former; on-device pitch and
azimuth estimation from one photo and the phone's pose; and a study of which obstacles the
ground view finds that satellite imagery does not.
\end{abstract}

\section{Introduction}
\section{Related work}
\section{Dataset}
\section{Benchmark}
\section{Baselines}
% generated by `roofsight leaderboard`; do not edit by hand
\begin{tabular}{lrrrrrrl}
\toprule
Model & Params & Mask AP & Small-obst. R & Boundary F & CPU ms & iPhone ms & Dataset \\
\midrule
\bottomrule
\end{tabular}

\section{Roof geometry from one photo}
\section{Ground view versus satellite: obstacle recall}
\section{Limitations}
\section{Conclusion}

\end{document}
